%-----------------------------------------------------------------------------------------------------------------------------------------------%
%	The MIT License (MIT)
%
%	Copyright (c) 2021 Jitin Nair
%
%	Permission is hereby granted, free of charge, to any person obtaining a copy
%	of this software and associated documentation files (the "Software"), to deal
%	in the Software without restriction, including without limitation the rights
%	to use, copy, modify, merge, publish, distribute, sublicense, and/or sell
%	copies of the Software, and to permit persons to whom the Software is
%	furnished to do so, subject to the following conditions:
%	
%	THE SOFTWARE IS PROVIDED "AS IS", WITHOUT WARRANTY OF ANY KIND, EXPRESS OR
%	IMPLIED, INCLUDING BUT NOT LIMITED TO THE WARRANTIES OF MERCHANTABILITY,
%	FITNESS FOR A PARTICULAR PURPOSE AND NONINFRINGEMENT. IN NO EVENT SHALL THE
%	AUTHORS OR COPYRIGHT HOLDERS BE LIABLE FOR ANY CLAIM, DAMAGES OR OTHER
%	LIABILITY, WHETHER IN AN ACTION OF CONTRACT, TORT OR OTHERWISE, ARISING FROM,
%	OUT OF OR IN CONNECTION WITH THE SOFTWARE OR THE USE OR OTHER DEALINGS IN
%	THE SOFTWARE.
%	
%
%-----------------------------------------------------------------------------------------------------------------------------------------------%

%----------------------------------------------------------------------------------------
%	DOCUMENT DEFINITION
%----------------------------------------------------------------------------------------

\documentclass[a4paper,11pt]{article}

%----------------------------------------------------------------------------------------
%	PACKAGES
%----------------------------------------------------------------------------------------
\usepackage{url}
\usepackage{parskip} 	

\RequirePackage{color}
\RequirePackage{graphicx}
\usepackage[usenames,dvipsnames]{xcolor}
\usepackage[scale=0.9]{geometry}

\usepackage{tabularx}
\usepackage{enumitem}

\newcolumntype{C}{>{\centering\arraybackslash}X} 

\usepackage{supertabular}
\usepackage{tabularx}
\newlength{\fullcollw}
\setlength{\fullcollw}{0.47\textwidth}

\usepackage{titlesec}				
\usepackage{multicol}
\usepackage{multirow}

\titleformat{\section}{\Large\scshape\raggedright}{}{0em}{}[\titlerule]
\titlespacing{\section}{0pt}{10pt}{10pt}

\usepackage[style=authoryear,sorting=ydnt, maxbibnames=2, maxcitenames=4, mincitenames=1]{biblatex}
\usepackage[unicode, draft=false]{hyperref}
\definecolor{linkcolour}{rgb}{0,0.2,0.6}
\hypersetup{colorlinks,breaklinks,urlcolor=linkcolour,linkcolor=linkcolour}
\addbibresource{papers.bib}
\setlength\bibitemsep{1em}

\usepackage{fontawesome5}

\newenvironment{jobshort}[2]
    {
    \begin{tabularx}{\linewidth}{@{}l X r@{}}
    \textbf{#1} & \hfill &  #2 \\[3.75pt]
    \end{tabularx}
    }
    {
    }

\newenvironment{joblong}[2]
    {
    \begin{tabularx}{\linewidth}{@{}l X r@{}}
    \textbf{#1} & \hfill &  #2 \\[3.75pt]
    \end{tabularx}
    \begin{minipage}[t]{\linewidth}
    \begin{itemize}[nosep,after=\strut, leftmargin=1em, itemsep=3pt,label=--]
    }
    {
    \end{itemize}
    \end{minipage}    
    }

%----------------------------------------------------------------------------------------
%	BEGIN DOCUMENT
%----------------------------------------------------------------------------------------
\begin{document}

\pagestyle{empty} 

\begin{tabularx}{\linewidth}{@{} C @{}}
\center{\footnotesize Last updated: \today} \hfill {\Huge{\textbf{Minh-Hai Nguyen}}} \\[7.5pt]
\href{https://www.linkedin.com/in/mh-nguyen712/}{\raisebox{-0.05\height}\faLinkedin\ mh-nguyen712} \ \(|\) \ 
\href{https://mh-nguyen712.github.io/}{\raisebox{-0.05\height}\faGlobe \ mh-nguyen712.github.io} \ \(|\) \ 
\href{mailto:nguyenhai7120qh@gmail.com}{\raisebox{-0.05\height}\faEnvelope \ nguyenhai7120qh@gmail.com} \ \(|\) \ 
\href{tel:+33 7 69 60 25 04}{\raisebox{-0.05\height}\faMobile \ +33 7 69 60 25 04} \\
\end{tabularx}

\vspace*{-0.3cm}
\section{Summary}
    I am a last year PhD candidate in the MAMBO team at Toulouse University (expected graduation: Dec. 2026), advised by \href{https://pierre-weiss.github.io/}{Prof. Pierre Weiss} and \href{https://edouardpauwels.fr/}{Prof. Edouard Pauwels}.
    My research focuses on computational imaging, mathematical analysis, and modern generative modeling (including score-based diffusion models and flow matching). I specialize in developing theoretical frameworks and deep learning models for inverse problems in imaging, with applications in microscopy and satellite imaging. I'm also interested in generative modeling for image generation, image reconstruction and vision-language models (VLMs). 

\section{Education}
\begin{tabularx}{\linewidth}{@{}l X@{}}	
2023 - present & Toulouse University, France \hfill \\
            & PhD in Applied Mathematics, advised by \href{https://pierre-weiss.github.io/}{Prof. Pierre Weiss} and \href{https://edouardpauwels.fr/}{Prof. Edouard Pauwels}. \\
2018 - 2023 & INSA Toulouse, France \hfill  \\
            & Master of Engineering in Applied Mathematics (incl. bachelor), top 3\% of the year \\
\end{tabularx}

\section{Work Experience}
\begin{jobshort}{Teaching assistant -- INSA Toulouse}{Dec 2023 - May 2026}
Courses: Signal Processing, Stochastic Optimization and Machine Learning
\end{jobshort}

\begin{jobshort}{Research Engineer -- \href{https://agenium-space.com/}{Agenium Space}}{Dec 2022 - Sept 2023}
Correcting micro-vibrations in satellite imaging with pushbroom cameras.
Modeled the physical imaging process, derived theoretical bounds for image restoration, and designed diffusion-based generative methods to improve real data from Sentinel-2B.
\end{jobshort}

\begin{jobshort}{Research Intern -- \href{https://www.math.univ-toulouse.fr/fr/}{Toulouse Mathematics Institute}}{Jun 2022 - Sept 2022}
Product-convolution neural networks: applications to image deblurring with space-varying blurs, with \href{https://pierre-weiss.github.io/}{Pierre Weiss} and \href{https://pescande.perso.math.cnrs.fr/}{Paul Escande}.
Formulated local quantitative stability properties and designed a deep learning-based framework for deblurring with complex spatially varying blurs.
\end{jobshort}


\section{Open-source Projects}
\begin{tabularx}{\linewidth}{ @{}l r@{} }
\textbf{Blind deblurring in Microscopy} & \hfill \href{https://deepblur.mambo-utoulouse.fr/}{Online demo}
\\[2pt]
\multicolumn{2}{@{}X@{}}{
	Designing mathematical and deep learning methods to jointly estimate the sharp image and the unknown blur kernel.
    Focused on rigorous physical modeling of optical systems, PSF estimation and image deblurring. 
    The framework supports various imaging modalities (RIM, TIRF, confocal) and space-varying blurs..
}  
\\ \\
\textbf{\href{https://deepinv.github.io/deepinv/}{DeepInverse} -- Core maintainer} & \hfill \href{https://github.com/deepinv/deepinv}{Link to Github} 
\\[2pt]
\multicolumn{2}{@{}X@{}}{
	An open-source library (\href{https://pytorch.org/blog/deepinverse-joins-pytorch-ecosystem/}{PyTorch Ecosystem}) for solving imaging inverse problems with state-of-the-art AI.
    Contributed advanced mathematical solvers, diffusion \& flow matching models, and algorithmic designs ensuring research reproducibility and robust theoretical baselines.
}  \\
\end{tabularx}

\section{Selected Publications}
\begin{refsection}[papers.bib]
\nocite{*}
\printbibliography[heading=none]
\end{refsection}

\section{Skills}
\textbf{Research:}
computational imaging, generative models (diffusions, flows, VLMs) and representation learning. \\
\textbf{Programming language:} Python, C++, SQL \\
\textbf{Deep learning frameworks:}
PyTorch, JAX, CUDA/Triton kernels, HuggingFace ecosystem, 
distributed training (DDP, FSDP), profiling (Nsight, torch-profiler), numpy, scikit-learn, pandas, HPC clusters (\href{http://www.idris.fr/eng/jean-zay/jean-zay-presentation-eng.html}{Jean-Zay}). \\
\textbf{Tools:}
Linux, Git/GitHub, Docker, open-source ecosystem. \\
\textbf{Languages:}
English (professional), French (fluent), Vietnamese (native).

\section{Talks and Presentations}
\begin{itemize}
    \item \href{https://icml.cc/}{ICML 2026 (Poster)} \hfill Seoul, South Korea, Jul. 2026
    \\ Title: An analytic theory of convolutional neural network inverse problems solvers
    \item \href{https://events.linuxfoundation.org/pytorch-conference-europe/program/schedule/}{PyTorch Conference Europe 2026 (Speaker)} \hfill Paris, France, Apr. 2026
    \\ Title: DeepInverse and Computational Imaging
    \item Blind inverse problems in imaging: from foundations to applications (Oral) \hfill Marseille, France, Oct 2025
    \\ Title: How diffusion prior landscapes shape the posterior in blind deconvolution
\end{itemize}

\section{Academic Service}
\subsection*{Reviewing}
\begin{itemize}
    \item Conferences: ICML 2026, CVPR 2025, GRETSI 2025
    \item Journals: Journal of Mathematical Imaging and Vision (JMIV)
\end{itemize}

\section{Distinctions}
\begin{itemize}
	\item French National Prize for Open Science Library for DeepInverse \hfill Nov 2024
	\item Doctoral Research Fellowship, French Ministry of Higher Education and Research \hfill 2023 - 2026
	\item IA Pau Winner: efficient language model (50+ teams) \hfill 2023 
\end{itemize}

\end{document}